\begin{blocksection}
\question Write \lstinline{count_paths}, which takes a tree \lstinline{t} of integers and an integer \lstinline{total}, and returns the number of root-to-leaf paths whose labels sum to \lstinline{total}.

\begin{lstlisting}
def count_paths(t: Tree[int], total: int) -> int:
    """
    >>> t = Tree(1, [Tree(3),
    ...              Tree(3, [Tree(4),
    ...                       Tree(0)]),
    ...              Tree(4, [Tree(1)])
    ...             ]
    ...     )
    >>> count_paths(t, 4)
    2
    >>> count_paths(t, 8)
    1
    >>> count_paths(t, 6)
    1
    >>> count_paths(t, 5)
    0
    """
    if ______________________________________:

        return ______________________________

    return ____________________________________________________________
\end{lstlisting}

\begin{solution}[1in]
\begin{lstlisting}
def count_paths(t: Tree[int], total: int) -> int:
    if is_leaf(t):
        return 1 if t.label == total else 0
    return sum([count_paths(b, total - t.label) for b in t.branches])
\end{lstlisting}
\end{solution}

\end{blocksection}

\begin{questionmeta}
\textbf{Teaching Tips}\par\nopagebreak[4]
Suggested Time: 10 min; Difficulty: Medium
\nopagebreak[4]
\begin{itemize}
    \item \textbf{Shrink the problem:} Once we use the root label, each branch only needs to produce a sum of \lstinline{total - t.label}.
    \item \textbf{Base case:} A leaf is a complete path; it counts exactly when its label equals the remaining total.
    \item \textbf{Combine:} Add the counts from all branches, as in \lstinline{count_leaves} from lecture.
    \item \textbf{Stopping early:} A path must end at a leaf. \lstinline{count_paths(t, 5)} is the only doctest that catches a solution that returns 1 as soon as the remaining total matches an internal node's label: the path 1, 4 sums to 5 but continues to the leaf 1. Ask students to trace it.
\end{itemize}
\end{questionmeta}
